\documentclass[12pt]{article}
\usepackage[margin=0.6in,top=0.85in,bottom=0.55in]{geometry}
\usepackage{lmodern}
\usepackage{microtype}
\usepackage{fancyhdr}
\usepackage{titlesec}
\usepackage{enumitem}
\usepackage{lastpage}
\usepackage[hidelinks]{hyperref}

\setlength{\parindent}{0pt}
\setlength{\parskip}{0.25em}
\setlength{\headheight}{26pt}
\titleformat{\section}{\large\bfseries\scshape}{}{0pt}{}
\pagestyle{fancy}
\fancyhf{}
\fancyhead[C]{\footnotesize Student Name: \hspace{2.3in} \quad Assignment ID: U1L1LE \quad Page \thepage\ of \pageref{LastPage}}
\renewcommand{\headrulewidth}{0.5pt}

\newcommand{\answerbox}[2]{% #1 = label, #2 = height
\vspace{0.12em}
\noindent\textbf{#1}\par
\vspace{0.04em}
\noindent\fbox{%
\begin{minipage}[t][#2]{\dimexpr\textwidth-2\fboxsep-2\fboxrule}
\rule{\textwidth}{0pt}
\vspace{#2}
\end{minipage}}
\par\vspace{0.22em}}

\begin{document}

\begin{center}
{\LARGE\bfseries Week 1 Lit Example Worksheet}\par\vspace{0.05em}
{\large\scshape Macbeth: Prophecy, Possibility, and Proof}\par\vspace{0.12em}
\rule{0.78\textwidth}{0.5pt}
\end{center}

\textbf{Purpose:} Apply Week 1's Whately method to Macbeth. Do not summarize the plot. For each answer, identify the logical movement: conclusion, stated reason, hidden assumption, and whether the conclusion follows.

\textbf{Directions:} Use the Lit Example Reader and quote briefly when helpful. Do not copy long blocks of Shakespeare. Write in complete, precise sentences. The strongest answers distinguish what a speaker \textit{wants}, \textit{fears}, or \textit{imagines} from what the speaker has actually proved.

\section*{Part 1: Prophecy and First Inference}

\textbf{Question 1.1}\\[-0.15em]
In Act 1, Scene 3, Macbeth has heard the witches' prophecy and then learns he is Thane of Cawdor. Identify one conclusion Macbeth is tempted to draw from this news, the stated reason that supports it, and the hidden assumption that would have to be true for the conclusion to follow.

\answerbox{ANSWER LABEL: Part 1 Q1}{2.15in}

\textbf{Question 1.2}\\[-0.15em]
Banquo warns that truthful predictions can still lead people toward harm. Reconstruct Banquo's warning as an argument: conclusion, stated reason, hidden assumption, and follow-test.

\answerbox{ANSWER LABEL: Part 1 Q2}{2.15in}

\newpage

\section*{Part 2: Macbeth as His Own Critic}

\textbf{Question 2.1}\\[-0.15em]
In Act 1, Scene 7, Macbeth gives reasons against murdering Duncan. Choose two reasons he gives and use them to reconstruct Macbeth's argument against the murder. Include conclusion, stated reasons, hidden assumption, and follow-test.

\answerbox{ANSWER LABEL: Part 2 Q1}{2.45in}

\textbf{Question 2.2}\\[-0.15em]
Macbeth says he has ``no spur'' except ``vaulting ambition.'' Explain the logical difference between a motive and a proof. Then apply that distinction to Macbeth's situation without telling the story back in summary form.

\answerbox{ANSWER LABEL: Part 2 Q2}{2.35in}

\newpage

\section*{Part 3: Testing What Actually Follows}

\textbf{Question 3.1}\\[-0.15em]
Choose one moment from either Macbeth passage where a speaker could move too quickly from possibility, fear, desire, or prediction to a larger conclusion. Name the possible conclusion, the actual support, the hidden assumption, and whether the conclusion follows.

\answerbox{ANSWER LABEL: Part 3 Q1}{2.55in}

\textbf{Question 3.2}\\[-0.15em]
Put Whately's central question to Macbeth: \textit{What follows from what?} Identify one claim in the two passages that is reasonably supported and one claim or possible claim that is not yet proved. Explain the difference logically.

\answerbox{ANSWER LABEL: Part 3 Q2}{2.45in}

\end{document}
